% Generated draft: requires mathematical and source review.
\documentclass[11pt]{article}
\usepackage[margin=25mm]{geometry}
\usepackage{fontspec}
\setmainfont{Latin Modern Roman}
\usepackage{amsmath,amssymb,mathtools,unicode-math}
\setmathfont{Latin Modern Math}
\usepackage{enumitem,needspace,xurl,hyperref}
\setlength{\parindent}{0pt}
\setlength{\parskip}{5pt}
\setlength{\emergencystretch}{3em}
\newcommand{\N}{\mathbb{N}}
\newcommand{\Z}{\mathbb{Z}}
\newcommand{\Q}{\mathbb{Q}}
\newcommand{\R}{\mathbb{R}}
\newcommand{\C}{\mathbb{C}}
\newcommand{\F}{\mathbb{F}}
\newcommand{\A}{\mathbb{A}}
\newcommand{\PP}{\mathbb{P}}
\newcommand{\E}{\mathbb{E}}
\newcommand{\eps}{\varepsilon}
\newcommand{\dd}{\,\mathrm{d}}
\begin{document}
\section*{UnsolvedMath ID: 4600048}
\textbf{Source problem code:} catalogue entry
\subsection{Definitions and mathematical statement}
A Salem number is a real algebraic integer $\beta>1$ whose conjugates other than $\beta$ consist of $\beta^{-1}$ and numbers on the unit circle, with at least one conjugate on that circle. Define the beta-transformation
\[
 T_\beta:[0,1)\to[0,1),\qquad
 T_\beta(x)=\beta x-\lfloor\beta x\rfloor.
\]
A periodic point means an $x$ with $T_\beta^k(x)=x$ for some integer $k\ge1$. This is distinct from eventual periodicity, which permits $T_\beta^{j+k}(x)=T_\beta^j(x)$ for some $j\ge0$ and $k\ge1$.

The literal assertion supplied in the record is
\[
 \forall\beta\text{ a Salem number},\qquad
 \bigcup_{k\ge1}\{x\in[0,1):T_\beta^k(x)=x\}
       =\Q\cap[0,1).
\]
This is retained as the extracted target, with a material consistency issue: a periodic-point set is invariant under $T_\beta$, whereas multiplication modulo one by irrational $\beta$ does not preserve $\Q\cap[0,1)$. Thus the literal equality requires checking against the original source before being accepted as the intended conjecture. Possible distinctions involving eventual rather than pure periodicity, the coefficient field of represented numbers, or a symbolic beta-shift formulation are not resolved or silently substituted.

The research summary links the question to Salem beta-shifts being sofic, to the orbit of $1$ being finite and $\beta$ being a Parry number, and also mentions shifts of finite type, ``restricted'' beta-shifts, ``finiteness,'' and an incomplete finite-orbit-complexity claim. It asserts equivalences among these formulations but supplies no exact definitions, hypotheses, or references verifying those equivalences. In particular the orbit of the endpoint $1$ needs a convention beyond the displayed half-open-domain transformation. Those claims cannot repair or replace the literal periodic-point equality.
\subsection{Short English statement}
The extracted statement asks whether the periodic points of every Salem beta-transformation are precisely the rational points of its interval. Its periodic-versus-eventual wording and rational-point target require source correction or confirmation because the stated equality has an invariance inconsistency.

The registry and research summary label the problem open and report no resolution through 2026, but the summary mixes several different beta-expansion and symbolic-dynamical properties. Its status assessment does not establish the literal equality or its asserted equivalence to those other questions. Provenance records no later authoritative resolution check of the 2008 source.
\subsection{Sources}
\begin{itemize}
\item Local UnsolvedMath registry, permanent ID 4600048, source code catalogue entry, ``Salem beta-transformations,'' proposed year recorded as 2008; no proposer is specified. Dataset curation and metadata: Ulam AI and UnsolvedMath Contributors, CC BY 4.0.
\item Mike Boyle, \emph{Open Problems in Symbolic Dynamics} (2008), Problem/Question/Conjecture 28.1, \url{https://www.math.umd.edu/~mboyle/papers/openfinalsub3nov2007.pdf}; pdf-text extraction, accessed 2026-07-29.
\item The supplied research summary gives no separate primary citation for its Salem-number, beta-expansion, soficity, or periodicity assertions. No missing reference or corrected conjecture is invented.
\item Source provenance marks status evidence NEEDS\_REVIEW and rights NEEDS\_REVIEW: the source presents the item as open, no later authoritative resolution check is recorded, and redistribution permission was not established despite verified public access. This qualification is separate from the dataset licence.
\end{itemize}
\end{document}
